\section{Distributed Market Clearing via Exchange ADMM}

To solve the stochastic coordination problem while preserving the privacy of residential FSPs, we adopt Exchange ADMM (Boyd et al.~\cite{stephen_distributed_2011}, §7.3.2) with $J = N+1$ agents: Agent~0 is the aggregator and agents $i \in \{1,\dots,N\}$ are the FSPs.

\subsection*{Centralized Problem and Coupling Constraint}

Let $s \in \mathcal{P} = \{\mathrm{AM},\,\mathrm{PM},\,\mathrm{AM{+}PM}\}$ index the GDP
plage scenarios with probabilities $\{\omega^s\}$, and let
\begin{equation}
  \pi_k \;=\; \sum_{\substack{s \in \mathcal{P} \\ k \in \mathcal{K}^s}} \omega^s
\end{equation}
denote the marginal activation probability of period $k \in \mathcal{K}$.
The centralized problem couples the aggregator and the FSPs through the \emph{delivery
constraint}: every unit committed must be supplied by FSPs or backed by CLC,
\begin{equation}
\label{eq:coupling}
\begin{aligned}
  \min_{\Xi_0,\,\{\Xi_i\}} \quad & f_0(\Xi_0) + \sum_{i=1}^{N} f_i(\Xi_i) \\[4pt]
  \text{s.t.} \quad
    & q_k^s \;=\; \sum_{i=1}^{N} F_{i,k} \;+\; c_k^s,
      \quad \forall\,s \in \mathcal{P},\;k \in \mathcal{K}^s
      \quad :\lambda_k
\end{aligned}
\end{equation}
where $q_k^s \leq \eta\cdot dt$ is the total delivery in scenario~$s$, $F_{i,k}$ is
FSP~$i$'s first-stage (non-anticipative) offer, and $c_k^s \geq 0$ is the CLC backup.
The dual variable $\lambda_k$ is the endogenous clearing price.
The Exchange ADMM decomposes \eqref{eq:coupling} by dualizing this constraint:
each agent solves a local subproblem augmented with a proximal penalty that drives
iterates towards market clearing.

\subsection*{ADMM Updates (Gauss–Seidel)}

The augmented Lagrangian introduces penalty $\rho > 0$.
Let $\mathbf{F}^n := \sum_i \mathbf{F}_i^n$ denote the aggregate FSP offer at iteration~$n$.
The Gauss–Seidel variant proceeds in three steps:

\paragraph{Step 1 — Aggregator.}
With $\mathbf{F}^{n}$ fixed, update $\bigl(\eta,\,c_{\max},\,\{q^s,\,c^s\}\bigr)$,
where $q_k^s := F_k^n + c_k^s$ is total delivery and
$r_k^s := \eta\cdot dt - q_k^s \geq 0$ is the shortfall (internal to $f_0$,
penalized as $p_{\mathrm{dev}}\cdot r_k^s$; determined by $\eta$ and $c^s$ once the FSP
signal $F_k^n$ is fixed):
\begin{equation}
\label{eq:agg_update}
  \{\Xi_0\}^{n+1} :=
    \arg\min_{\Xi_0}
    \left(
      -f_0 \;+\;
      \sum_{s \in \mathcal{P}}\omega^s\sum_{k \in \mathcal{K}^s}
      \left[
        \lambda_k^n q_k^s
        \;+\;
        \frac{\rho}{2}\,\bigl\|q_k^s - \sigma_k^{0}\bigr\|^2
      \right]
    \right)
\end{equation}
Since $r_k^s = \eta\cdot dt - q_k^s = \eta\cdot dt - F_k^n - c_k^s$,
the proximal target depends only on $\eta^n$ and $F_k^n$:
\begin{equation}
  \sigma_k^{0} \;=\; \frac{J-1}{J}\,\pi_k\!\left(\eta^n\cdot dt - F_k^{n}\right)
                \;-\; \frac{1}{J} F_k^{n},
  \qquad k \in \mathcal{K}.
\end{equation}

\paragraph{Step 2 — FSPs (parallel).}
Define the Gauss–Seidel correction using the freshly updated aggregator:
\begin{equation}
\label{eq:xbar_half}
  \delta_k \;:=\; \frac{1}{J}\!\left[
    \pi_k\!\left(\eta^{n+1}\cdot dt - F_k^{n}\right) - F_k^{n}
  \right].
\end{equation}
Each FSP~$i$ solves independently:
\begin{equation}
\label{eq:fsp_update}
  \{\Xi_i\}^{n+1} :=
    \arg\min_{\Xi_i}
    \left(
      f_i
      \;-\; \sum_{k \in \mathcal{K}} \pi_k\lambda_k^n F_{i,k}
      \;+\; \frac{\rho}{2}\sum_{s}\omega^s
            \bigl\|\tilde{\mathbf{F}}_{i}^s - \boldsymbol{\sigma}_{i}\bigr\|^2
    \right)
\end{equation}
with proximal target $\sigma_{i,k} = F_{i,k}^{n} + \delta_k$,
subject to non-anticipativity $\tilde{F}_{i,k}^s = F_{i,k}\ \forall s \in \mathcal{P}$.

\paragraph{Step 3 — Dual update.}
\begin{equation}
\label{eq:dual_update}
  \lambda_k^{n+1} \;=\; \lambda_k^n \;+\; \frac{\rho\,\pi_k}{J}\,r_{p,k}^{n+1},
  \qquad
  r_{p,k}^{n+1} \;:=\; \mathbb{E}_s\!\left[q_k^{s,n+1} - c_k^{s,n+1}\right]
                        \;-\; \pi_k F_k^{n+1}
                 \;=\; \pi_k\!\left(F_k^n - F_k^{n+1}\right),
\end{equation}
where the last equality holds because the aggregator uses $F_k^n$ as a fixed parameter,
so $q_k^{s,n+1} - c_k^{s,n+1} = F_k^n$ identically.

\subsection*{Convergence}

Iteration stops when both residuals fall below their tolerances $\varepsilon^{\mathrm{pri}}$ and
$\varepsilon^{\mathrm{dual}}$:
\begin{equation}
\label{eq:convergence}
  \|\mathbf{r}_p^n\|_2 \;\leq\; \varepsilon^{\mathrm{pri}},
  \qquad
  \|\mathbf{r}_d^n\|_2 \;\leq\; \varepsilon^{\mathrm{dual}},
\end{equation}
where
\begin{align}
  r_{p,k}^n &\;=\; \pi_k\!\left(F_k^{n-1} - F_k^{n}\right)
             && \text{(coupling constraint violation)}, \\[2pt]
  r_{d,k}^n &\;=\; \rho\!\left(F_k^n - F_k^{n-1}\right)
             && \text{(dual: FSP offer change)}.
\end{align}
The penalty $\rho$ is adapted at each iteration following Boyd §3.4.1:
scaled by $\tau$ when $\|\mathbf{r}_p\| > \mu\|\mathbf{r}_d\|$, and divided when
$\|\mathbf{r}_d\| > \mu\|\mathbf{r}_p\|$ (parameters: $\tau = 1.2$, $\mu = 20$).

Since FSP constraints are coupled only through $\lambda_k$, convergence yields a
Generalised Nash Equilibrium in which each agent's local decisions align with the
global market commitment while FSP cost data remain private.
